\scanpage{201}
\hfill -- \mspage{57} -- \hfill IX

5.2.6. Supposons maintenant que $A_K$ ait une réduction semi-stable sur $S$, de sorte que nous pouvons prendre pour $\ell$ n'importe quel nombre premier (distinct ou non de $p$). D'ailleurs, en vertu du critère galoisien de réduction semi-stable 3.5, on a pour $\ell\neq p$
\[
\text{(5.2.7)}\qquad (T_\ell(A'_K)^f)^\perp\quad (=T_\ell(A_K)^t)\subset T_\ell(A_K)^f\quad .
\]
Il reste à prouver, pour terminer la preuve de 5.2, que cette inclusion est valable sans restriction sur $\ell$.

En effet, le membre de gauche est canoniquement isomorphe à $D(T_\ell(A'_K)/T_\ell(A'_K)^f)$, et est donc de ``rang'' (ou ``hauteur'') égal à $2n-2\alpha-\mu=\mu$ (notations de (2.1.9), où par hypothèse $\lambda=0$), qui est aussi celui de $T_\ell(A_K)^t$. Comme l'inclusion
\[
\text{(*)}\qquad T_\ell(A_K)^t\subset T_\ell(A_K)
\]
est \uline{stricte}, dans le sens qu'elle induit déjà un monomorphisme pour tous les termes du systèmes projectifs envisagés, il s'ensuit aisément que le premier membre de (*) est maximal parmi les sous-pro-groupes de Barsotti-Tate de $T_\ell(A_K)$ ayant même rang que lui, de sorte qu'on a bien l'égalité de (5.2.7), qui équivaut à l'inclusion entre les termes extrêmes de 5.2.7 par le théorème d'orthogonalité.

5.3. Utilisant la correspondance par orthogonalité entre sous-pro-groupes de Barsotti-Tate stricts (au sens précisé plus haut) des pro-groupes duaux $T_\ell(A_K)$ et $T_\ell(A'_K)$, on déduit formellement de (5.2.2) que l'accouplement (1.0.2) entre ces pro-groupes induit un accouplement
\[
\text{(5.3.2)}\qquad T_\ell(A_K)^{ab}\times T_\ell(A'_K)^{ab}\longrightarrow T_\ell(\mathbf G_{mK})
\]
\margin{375}
\margin{375}
\margin{200}
